\section*{Abstract}

Disposable vapes have become a hazardous form of urban micro-litter in the United Kingdom. Although small and visually unobtrusive, they contain lithium-ion batteries, metals, plastics, electronic components and e-liquid residues, meaning that missed or incorrectly discarded devices can create environmental and waste-processing risks. This dissertation investigates how disposable vapes can be identified and operationalised as hazardous micro-WEEE in UAV-oriented imagery for urban environments.

The work is structured around three connected research workstreams. First, a Dataset Construction Engine (DCE) was developed to support image scraping, synthetic data creation, upload, curation, SAM2-assisted annotation, metadata enrichment, audit diagnostics and COCO/YOLO export. The DCE was designed to address the practical bottleneck that niche computer-vision datasets often do not already exist and must be constructed before model evaluation can begin. Second, the DCE was used to construct UAVape, a two-class UAV-oriented vape/lighter litter dataset containing 1,204 real images, 1,328 vape boxes and 712 lighter boxes, alongside separate scraped and synthetic training-source pools. A leakage-aware split and unified YOLO benchmark were then used to evaluate model family, model scale, input resolution, tiled inference and data-source ablations. The selected YOLO26s protocol achieved a held-out vape AP50 of 0.882 and vape AP50--95 of 0.685 on the sealed real-image test set. Third, a cleaner-facing workflow prototype, the \textit{UAVape Cleanup Planner}, explored how detections could be translated into evidence crops, map points, hazard labels, task assignment and field status updates.

The results show that vape detection is not only a model-selection problem, but a dataset-construction and workflow-translation problem. The main contribution is therefore an end-to-end methodology for constructing, evaluating and operationalising a new hazardous micro-litter class, while remaining bounded by the study's limitations: the detector was evaluated offline on static imagery, and the workflow prototype was assessed through stakeholder walkthrough feedback rather than live municipal deployment.
